\documentclass[10pt,a4paper]{article}
\usepackage[T1]{fontenc}
\usepackage[utf8]{inputenc}
\usepackage[margin=24mm]{geometry}
\usepackage{amsmath,amssymb,amsthm,booktabs}
\usepackage{hyperref}
\hypersetup{colorlinks=true,linkcolor=blue,citecolor=blue,urlcolor=blue,
  pdftitle={Exact Egyptian-fraction certificates and overlap-constrained search},
  pdfauthor={Kyunghan Bae}}
\setlength{\parindent}{0pt}
\setlength{\parskip}{5pt}
\setlength{\emergencystretch}{2em}
\newtheorem{proposition}{Proposition}
\newtheorem{lemma}[proposition]{Lemma}
\newcommand{\card}[1]{\lvert #1\rvert}
\newcommand{\code}[1]{\texttt{\detokenize{#1}}}
\title{Exact Egyptian-fraction certificates\\and overlap-constrained search}
\author{Kyunghan Bae}
\date{12 October 2026\\[3pt]
  \small DOI: \href{https://doi.org/10.5281/zenodo.23305279}{10.5281/zenodo.23305279}}

\begin{document}
\maketitle
\begin{abstract}
We give explicit expansions of $1$ into distinct unit fractions establishing
$k(18)\leq34$, $k(19)\leq36$, $k(20)\leq38$ and $k(21)\leq39$.
Each upper bound has a finite certificate verified by exact rational and integer
arithmetic. An elementary cancellation lemma converts complete families of local
shortening exclusions into overlap restrictions on every shorter representation.
For the $34$-term reference, complete five-to-four and fixed-core six-to-five
searches give capacities $28$ and $27$. Classical quota optimization supplies
exact reciprocal upper envelopes. A fixed $53$-branch pilot compares three modes
of one C++/GMP engine with a pinned original Sage implementation: the matched modes
exhaust all $53$ branches; Sage exhausts $48$ and times out on five, with no
contradictory completed answers. The existence of a $33$-term expansion at $N=18$
and the asymptotic question in Erd\H{o}s problem~295 remain open.
\end{abstract}

\section{Problem and contribution}
For a positive integer $N$, let $k(N)$ be the least cardinality of a finite set
$D$ of distinct integers at least $N$ satisfying
\[
  \sum_{d\in D}\frac1d=1.
\]
The main results are four explicit finite identities. Their term counts compare
with the identified computational certificates as follows:
\begin{center}
\begin{tabular}{rrr}
\toprule
$N$ & Cited certificate term count & Term count here\\
\midrule
18&35&34\\19&38&36\\20&39&38\\21&40&39\\
\bottomrule
\end{tabular}
\end{center}
The $N=18$ construction shortens the $35$-term expansion recorded by White with
Claude~\cite{white}. The comparisons for $N=19,20,21$ use the vectors published by
grind-35 on 24 September 2026~\cite{grind}. All source vectors used in these
comparisons were rechecked with exact rational arithmetic. Each resulting
expansion includes $N$ itself, so it also gives an upper bound for OEIS A192881,
whose definition requires the first denominator to equal $N$. The comparison is
to the specified records.

Fixed-length enumeration and divisor-based two-term completion are
classical~\cite{eppstein}. Alekseyev provides related exact enumeration software
in Sage~\cite{aleks}. Greedy optimization on matroids and laminar capacity
families is established theory~\cite{edmonds,fife}; membership-pattern dynamic
programming and weak duality are standard methods. Here they accompany the
explicit certificates, complete local-exclusion records and their elementary
transfer into two concrete overlap restrictions. Finite allowed-denominator
enumeration~\cite{martin} is a related problem; the representation searches here
allow unbounded denominators.

\section{Four finite certificates}
\begin{proposition}\label{prop:certificates}
The following sets $D_N$ have the stated cardinalities, consist of distinct
integers at least $N$, and satisfy $\sum_{d\in D_N}1/d=1$.
\end{proposition}

\textbf{$N=18$, $34$ terms:}
\begin{quote}\small
\begin{verbatim}
18,19,20,21,22,23,24,25,26,27,28,29,30,31,
32,33,34,35,36,38,39,40,43,44,45,46,
48,50,52,54,63,275,14143068,10457416992
\end{verbatim}
\end{quote}
\textbf{$N=19$, $36$ terms:}
\begin{quote}\small
\begin{verbatim}
19,20,21,22,23,24,25,26,27,28,29,30,31,32,
33,34,35,36,38,39,40,42,43,44,45,46,48,50,
52,54,55,66,70,275,14143068,10457416992
\end{verbatim}
\end{quote}
\textbf{$N=20$, $38$ terms:}
\begin{quote}\small
\begin{verbatim}
20,21,22,23,24,25,26,27,28,29,30,31,32,33,
34,35,36,37,38,39,40,41,42,43,44,45,46,
48,50,51,52,53,54,72,9512,316569706,
745735753538714400,755732875917998385
\end{verbatim}
\end{quote}
\textbf{$N=21$, $39$ terms:}
\begin{quote}\small
\begin{verbatim}
21,22,23,24,25,26,27,28,29,30,31,32,33,34,
35,36,37,38,39,40,41,42,44,45,46,48,49,50,
52,53,54,55,56,57,58,64,3162192,
13244735025984,60692489317197426150
\end{verbatim}
\end{quote}
\begin{proof}
Let $L_N$ be the least common multiple of the displayed denominators. Exact
integer addition gives the following certificates:
\begin{center}
\begin{tabular}{rl}
\toprule
$N$ & $L_N=\sum_{d\in D_N}L_N/d$\\
\midrule
18 & \texttt{6209352774424800}\\
19 & \texttt{6209352774424800}\\
20 & \texttt{2772929439982611617430667744800}\\
21 & \texttt{207311948389939355504942400}\\
\bottomrule
\end{tabular}
\end{center}
Each denominator divides its stated $L_N$, and the integer quotients sum to
$L_N$. Division by $L_N$ proves the reciprocal identity. Cardinality, increasing
order and minimum denominator are checked directly from the lists.
\end{proof}
The authoritative lists and all these integers are stored in
\code{certificates-extension.json}, with large integers encoded as decimal
strings. Verification of Proposition~\ref{prop:certificates} is independent of
the search that found the identities.

\section{Construction and classical completion}
For $N=18$, retain $30$ terms of the earlier $35$-term certificate~\cite{white}
and substitute
\begin{align*}
 &\frac1{42}+\frac1{640}+\frac1{337280}
     +\frac1{1735650}+\frac1{68479214496}\\
 &\hspace{15mm}=\frac1{46}+\frac1{275}
     +\frac1{14143068}+\frac1{10457416992}.
\end{align*}
Multiplication by $636856694812800$ gives $16160598373892$ on each side. The
four new denominators are distinct and absent from the retained terms. For
$N=19$, replace $18$ and $63$ in $D_{18}$ by $42,55,66,70$, using
\[
 \frac1{18}+\frac1{63}=\frac1{14}
 =\frac1{42}+\frac1{55}+\frac1{66}+\frac1{70}.
\]
For $N=20$ and $N=21$, retain respectively the first $36$ and $37$ denominators
of the cited $39$- and $40$-term certificates~\cite{grind}. Replace their remaining
three terms by the two-term completions
\begin{align*}
 \frac{151817}{56984589375772121632800}
 &=\frac1{745735753538714400}+\frac1{755732875917998385},\\
 \frac{35558057}{470956940228352916800}
 &=\frac1{13244735025984}+\frac1{60692489317197426150}.
\end{align*}

The classical completion rule~\cite{eppstein} follows from
\begin{equation}\label{eq:factor}
 (ax-b)(ay-b)=b^2.
\end{equation}
For a positive reduced residual $a/b$ and a lower endpoint $D$ exceeding the
retained denominators, a completion $a/b=1/x+1/y$, $D\leq x<y$, exists exactly
when a positive divisor $q$ of $b^2$ satisfies
\[
 \max(1,aD-b)\leq q<b,\qquad q\equiv-b\pmod a.
\]
Then $x=(b+q)/a$ and $y=(b+b^2/q)/a$. A solution gives $q=ax-b$ by
\eqref{eq:factor}. Conversely, coprimality and the congruence imply that both
denominators are integers; the inequalities give $D\leq x<y$. The case $a=1$
is immediate. Complete divisor enumeration decides completion for the specified
prefix, rather than all prefixes. The reported $N=20$ pair minimizes the larger
denominator among completions of its fixed retained prefix, not globally.

\section{An elementary transfer lemma}
Let $U$ be a universe, $w:U\to\mathbb R$ any weight function, and
$W(S)=\sum_{x\in S}w(x)$ for finite $S$. Fix $A\subseteq U$ with $\card A=m$,
integers $1\leq d\leq r\leq m$, and $C\subseteq A$ with $\card C\leq r$.
A target is a set $B$ with $\card B=m-d$ and $W(B)=W(A)$.

For $R\subseteq A$, $\card R=r$, an admissible replacement is a set $T$ such that
\[
 \card T=r-d,\qquad T\cap(A\setminus R)=\varnothing,
 \qquad W(T)=W(R).
\]
\emph{Removed elements may reappear in $T$}; only retained elements are
forbidden. This convention is essential when padding the difference between
two representations.

\begin{lemma}[Fixed-core transfer]\label{lem:transfer}
Suppose every $r$-element subset $R$ of $A$ containing $C$ has no admissible
replacement. Then every target $B$ satisfies
\[
 \card{B\cap(A\setminus C)}\leq m-r-1.
\]
\end{lemma}
\begin{proof}
For a fixed $R$, a replacement produces the target
$B=(A\setminus R)\cup T$ containing $A\setminus R$. Disjointness gives the
cardinality and cancellation gives the weight. Conversely, a target containing
$A\setminus R$ gives $T=B\setminus(A\setminus R)$ of cardinality $r-d$,
disjoint from the retained set and of weight $W(R)$. Thus replacement existence
is equivalent to the existence of a target containing the retained set.

Let $K=A\setminus B$. If $\card{K\cup C}\leq r$, extend $K\cup C$ to an
$r$-element subset $R\subseteq A$. Then $A\setminus R\subseteq B$, giving a
replacement by the preceding equivalence, contrary to the premise. Hence
$\card{K\cup C}\geq r+1$, and complements inside $A$ give
\[
 \card{B\cap(A\setminus C)}=m-\card{K\cup C}\leq m-r-1.
\]
Neither positivity nor rationality of the weights is used.
\end{proof}
For $C=\varnothing$, every target overlaps $A$ in at most $m-r-1$ elements.
This is an elementary cancellation lemma. The substantive computational premise
is the complete family of exclusions.

\section{Completed premises and reciprocal bounds}
Take $A=D_{18}$ and $C=\{14143068,10457416992\}$. The local searches use
distinct denominators at least $18$ and forbid exactly the retained terms:
\begin{center}\small
\begin{tabular}{llrr}
\toprule
Removed/added & Removed-subset family & Complete subsets & Seconds\\
\midrule
$5/4$ & Every $5$-subset of $A$ & $\binom{34}{5}=278256$ &95.781\\
$6/5$ & Every $6$-subset containing $C$ & $\binom{32}{4}=35960$ &257.625\\
\bottomrule
\end{tabular}
\end{center}
These families contain $314216$ completed subsets. In the second family, the
first log completes indices $0$--$33953$ and interrupts index $33954$. The resumed
run replays $33954$ and completes through $35959$. The interrupted attempt is not
counted as complete; the second time includes replay.
\code{search-records/verify_local_premises.py} audits this coverage and boundary.

There is no denominator cap. At a positive reduced remainder $a/b$ with $t\geq2$
terms left and lower endpoint $D$, the next denominator lies in the finite
necessary interval
\[
 \max(D,\lfloor b/a\rfloor+1)\leq x\leq\lfloor tb/a\rfloor.
\]
Exact harmonic bounds and proved arithmetic conditions prune branches. With two
terms left, $y=bx/(ax-b)$ is checked exactly, without an imposed maximum on the
final denominator. The full recurrence and arithmetic details are supplied in
\code{proofs.md} and \code{SEARCH_STATUS.md}.

Lemma~\ref{lem:transfer} gives, for every $33$-term representation $B$ with
denominators at least $18$,
\begin{equation}\label{eq:quotas}
 \card{B\cap A}\leq28,\qquad
 \card{B\cap(A\setminus C)}\leq27.
\end{equation}
These restrictions apply to the entire unbounded denominator universe. After a
prefix $P$, subtract its selected overlap counts from the capacities. For $t$
remaining denominators at least $D=\max(P)+1$, accept denominators in increasing
order when both residual capacities permit them. Their reciprocal sum is the
exact maximum under these nested quotas, by classical greedy optimization for a
truncated laminar matroid~\cite{edmonds,fife}. A negative residual capacity, or a
target above the maximum, excludes the branch.

For crossing finite quota groups, partition denominators by membership pattern.
Within each pattern and chosen count, use the smallest denominators; dynamic
programming over counts gives the exact upper envelope. Only the first $t$
members of each pattern are needed. This finite support evaluates a bound and
does not restrict the representation search. Nonnegative exact multipliers give
standard weak-duality bounds. Proofs, resource guards and executable checks are
supplied in \code{STRUCTURE.md} and \code{structure/}.

\section{A controlled original-Sage pilot}
The frozen corpus has $224$ prefixes, each extended toward a $33$-term
representation of $1$ with minimum denominator $18$:
\begin{itemize}
\item $190$ broad prefixes
$\{18,\ldots,28\}\cup\{30,\ldots,44\}\cup\{u,v\}$,
$45\leq u<v\leq64$;
\item $31$ structural prefixes: the first $27$ elements of $A$, followed by
$z=49,\ldots,79$;
\item three controls $\{18,\ldots,28\}\cup\{z\}$ for $z=31,32,33$.
\end{itemize}
The pilot includes all structural prefixes and controls, plus every tenth broad
prefix in lexicographic $(u,v)$ order, starting at ordinal zero: $53$ cases fixed
before timing.

The pilot used one Linux runner CPU and one repetition per case. The unchanged
original Sage source~\cite{aleks} was fetched at commit
\code{8680599e77683a9057ee55871a0893bb319300e8}, with the recorded SHA-256
checked before execution. It ran in SageMath $10.8$ with one map-reduce worker.
The wrapper called \code{one_rep} on each exact residual with
\code{distinct=True}, the correct lower endpoint and no maximum denominator.
Six API controls passed, including the exact set of all six four-term distinct
expansions of $1$ with denominators at least $2$.

One C++/GMP source ran three modes: OFF, full-reference capacity $28$, and nested
capacities $28/27$. Arithmetic, traversal, prime/quadratic pruning and memo limits
were shared. Ordinary cases had a five-second solver budget and the frontier
controls had $30$ seconds:
\begin{center}
\begin{tabular}{lrrr}
\toprule
Solver/mode & Exhausted & Timed out & Errors\\
\midrule
Original Sage&48&5&0\\
Matched GMP, OFF&53&0&0\\
Matched GMP, single quota&53&0&0\\
Matched GMP, nested quotas&53&0&0\\
\bottomrule
\end{tabular}
\end{center}
All $48$ branches completed by both programs agreed on absence. No contradictory
completed answer occurred. The five Sage timeouts remain censored baseline runs.
The completed matched searches exclude these specified branches, rather than
all $33$-term possibilities. Saved records, source/image hashes and the completed
workflow are available in~\cite{cloud}.

The inter-program comparison includes C++/GMP versus Sage/Python execution,
different tail routines, harmonic/prime/quadratic pruning and worker-management
costs. Its timing difference does not isolate the effect of overlap quotas.
The three matched modes provide that ablation. For example,
\code{structural-052} visits $9444$, $7276$ and $1$ states in OFF, single and
nested modes, respectively, with the same exhausted verdict. This demonstrates
a concrete quota exclusion in the selected family. A single pilot and selected
structural branches do not establish a general runtime advantage.

Solver timers exclude compilation, Sage startup/source loading and deriving the
local-exclusion premises. Sage API timing includes worker setup and cleanup;
separate wall times are retained. The historical premise times were measured
separately and are not hardware-normalized to the cloud pilot. Timeouts and late
answers remain in the records.

Independent unpruned rational enumeration checked $972$ conditional residual
states with one, two or three terms left: $349$ admitted a witness and $623$ were
exhausted. Verdicts agreed with the matched engine, and every returned witness
passed exact sum, order, distinctness and quota checks. These are implementation
and boundary controls, not a global $33$-term exclusion.

\subsection{Repeated local quota ablation}
A separate Windows run executed all $224$ cases in all three matched modes,
with three repetitions: $2016$ solver runs. The declared budget was two seconds
per ordinary case and $30$ seconds per control. Strict post-processing censored
nine raw timeouts and four answers returned after their deadline; all $2016$
records are retained. There were no errors or contradictory completed verdicts.
Exactly $213$ cases finished within budget in all nine runs.

For each case and mode, take the median of its three solver times. The table
reports the sum of these medians for OFF divided by the sum for nested quotas,
restricted in each row to cases whose nine runs all completed within budget.
A ratio above $1$ favors nested quotas on this conditional measure.
\begin{center}
\begin{tabular}{lrr}
\toprule
Family & Complete / declared & OFF / nested \\
\midrule
All & $213/224$ & $1.05244$ \\
Broad & $181/190$ & $1.01684$ \\
Structural & $29/31$ & $1.67216$ \\
Controls & $3/3$ & $0.92689$ \\
\bottomrule
\end{tabular}
\end{center}
The overall difference is modest, with the larger effect concentrated in the
selected structural family. The broad family changes little, while controls
take longer with nested quotas on this measure. These ratios concern online
solver time on a common completed subset; they exclude premise computation and
compilation, and do not measure the full suite including censored runs. For the same $213$ instances, sums of median external wall times were
$108.386984$ seconds for OFF and $109.278968$ for nested quotas. Process launch,
output collection and validation therefore removed the small aggregate
internal-time gain in this execution. The Windows repetitions and the Linux
Sage pilot are separate experiments.

\section{Reproducibility and remaining questions}
Code and records are available in~\cite{repo}. Python $3.9+$ and its standard
library suffice for \code{python verify-extension.py}; a separately written
BigInt check runs as \code{node verify-extension-bigint.mjs}.
The verifier checks rational sums, integer LCM identities, cardinality, order,
minima and source-to-result substitutions. Twenty deliberately corrupted records
are rejected. The general transfer and envelope checks run as
\code{python structure/general_theorem_oracle.py}. Local-exclusion audits and
search programs are preserved separately; search is unnecessary to verify
Proposition~\ref{prop:certificates}.

The accompanying release reproduces the previously reported exclusion of
exactly $32$ terms at $N=18$, using the same mathematical recurrence in separately
written exact-arithmetic implementations. The logged run has $814918$ visited
states; all two-term states are rejected before candidate enumeration.
Excluding $32$ terms also excludes every shorter expansion: repeatedly replace
its largest denominator $d$ by $d+1$ and $d(d+1)$, using
\[
 \frac1d=\frac1{d+1}+\frac1{d(d+1)},
\]
until $32$ terms are reached. The new denominators exceed every retained one.
Thus the recorded bounds are $33\leq k(18)\leq34$. Whether $k(18)$ is $33$ or
$34$ remains unresolved here.

Erd\H{o}s problem~295 asks whether
$k(N)-(e-1)N\to+\infty$~\cite{erdos}. These finite certificates and overlap
exclusions do not settle that question. Further progress would require a
completed minimum proof, a quantitatively stronger family of constructions or a
broader controlled study of the pruning rules.

\paragraph{Computational contribution.}
Kyunghan Bae initiated the project; OpenAI Codex agents implemented and ran the
searches, found the certificates, and assisted with verification and exposition.

\begin{thebibliography}{10}\small
\bibitem{erdos} Thomas Bloom, ed.
\emph{Erd\H{o}s Problem \#295}.
\url{https://www.erdosproblems.com/295}.

\bibitem{white} Patrick White with Claude.
\emph{Erd\H{o}s \#295 computational report}, including the earlier $35$-term
$N=18$ certificate.
\url{https://erdosproblemaday.com/day/295-k17-exact}.
The original construction and provenance are preserved in \code{NOTE.md}.

\bibitem{grind} grind-35.
\emph{Erdos 295 upper bounds N=11..22}, 24 September 2026.
\url{https://botnet.com/artifacts/c1db3093-569e-4e89-bced-fd38231c0268}.

\bibitem{eppstein} David Eppstein.
\emph{Algorithms for Egyptian Fractions: Brute Force Methods}.
\url{https://ics.uci.edu/~eppstein/numth/egypt/force.html}.

\bibitem{aleks} Max A. Alekseyev.
\emph{egypt\_rep.sage}, version 20251028,
commit \code{8680599e77683a9057ee55871a0893bb319300e8},
accessed 12 October 2026.
\url{https://github.com/maxale/egyptian-fractions/blob/8680599e77683a9057ee55871a0893bb319300e8/egypt_rep.sage}.

\bibitem{edmonds} Jack Edmonds.
Matroids and the greedy algorithm.
\emph{Mathematical Programming} \textbf{1} (1971), 127--136.
\url{https://doi.org/10.1007/BF01584082}.

\bibitem{fife} Tara Fife and James Oxley.
\emph{Laminar Matroids}, arXiv:1606.08354 (2016).
\url{https://arxiv.org/abs/1606.08354}.

\bibitem{martin} Gregory Martin and Yue Shi.
An algorithm for Egyptian fraction representations with restricted denominators.
\emph{Involve} \textbf{18} (2025), 1--23.
\url{https://doi.org/10.2140/involve.2025.18.1}.

\bibitem{cloud} \emph{Original Sage and matched-engine cloud pilot},
workflow 38176293521, with \code{cloud-comparison.json} and
\code{matched-oracle-validation.json}.
\url{https://github.com/qorud02/erdos-295-k18/actions/runs/38176293521}.

\bibitem{repo} Kyunghan Bae.
\emph{Egyptian-fraction certificates, code and records}.
\url{https://github.com/qorud02/erdos-295-k18}.
Original $N=18$ archive:
\url{https://doi.org/10.5281/zenodo.23297810}.
\end{thebibliography}
\end{document}
